\documentclass[letterpaper,11pt]{article}

\usepackage[T1]{fontenc}
\usepackage[sc,osf]{mathpazo}
\usepackage{geometry}
\usepackage[dvipdfmx]{hyperref}
\usepackage{xcolor}
\usepackage{titlesec}
\usepackage{tabularx}
\usepackage{enumitem}
\usepackage{parskip}
\usepackage{fancyhdr}

\def\name{Kento Kawaharazuka}
\def\footerlink{https://haraduka.github.io}

\definecolor{accent}{HTML}{1D4E89}
\definecolor{textgray}{HTML}{444444}
\definecolor{lightgray}{HTML}{E8EDF3}

\geometry{
  left=0.85in,
  right=0.85in,
  top=0.8in,
  bottom=0.85in
}

\hypersetup{
  colorlinks = true,
  urlcolor = accent,
  linkcolor = accent,
  pdfauthor = {\name},
  pdfkeywords = {humanoid robotics, tendon-driven robots, machine learning},
  pdftitle = {\name: Curriculum Vitae},
  pdfsubject = {Curriculum Vitae},
  pdfpagemode = UseNone
}

\pagestyle{fancy}
\fancyhf{}
\renewcommand{\headrulewidth}{0pt}
\fancyfoot[C]{\footnotesize Last updated: \today \quad|\quad \href{\footerlink}{\texttt{\footerlink}}}

\titleformat{\section}
  {\Large\bfseries\color{accent}}
  {}
  {0em}
  {}[\vspace{0.15em}\titlerule\vspace{0.35em}]

\titleformat{\subsection}
  {\normalsize\bfseries\color{accent}}
  {}
  {0em}
  {}

\setlist[itemize]{leftmargin=1.4em,itemsep=0.2em,topsep=0.2em}
\setlength{\parindent}{0pt}
\setlength{\parskip}{0.35em}

\newcommand{\cventry}[2]{%
  \noindent\begin{tabularx}{\linewidth}{@{}p{0.2\linewidth}X@{}}%
    \textcolor{textgray}{#1} & #2%
  \end{tabularx}\vspace{0.2em}%
}

\begin{document}

{\Huge\bfseries \name}\\[-0.1em]
{\large Lecturer (Junior Associate Professor)}\\
{\normalsize Next Generation Artificial Intelligence Research Center (AI Center) \\ Department of Mechano-Informatics, Graduate School of Information Science and Technology, The University of Tokyo}

\vspace{0.3em}
\color{lightgray}\rule{\linewidth}{1.4pt}\color{black}

\begin{tabularx}{\linewidth}{@{}X X@{}}
\textbf{Email:} \href{mailto:kawaharazuka@jsk.imi.i.u-tokyo.ac.jp}{kawaharazuka@jsk.imi.i.u-tokyo.ac.jp} & \textbf{Website:} \href{https://haraduka.github.io}{https://haraduka.github.io} \\
\textbf{Google Scholar:} \href{https://scholar.google.co.jp/citations?user=E75YHyUAAAAJ}{profile link} & \textbf{GitHub:} \href{https://github.com/haraduka}{github.com/haraduka}
\end{tabularx}

\section*{Employment}
\cventry{2025.10--Present}{\textbf{Visiting Researcher}, Robot System Team, RIKEN Center for Advanced Intelligence Project (RIKEN AIP), Japan}
\cventry{2025.02--Present}{\textbf{Lecturer (Junior Associate Professor)}, UTokyo AI Center and Department of Mechano-Informatics, Graduate School of Information Science and Technology, The University of Tokyo, Japan}
\cventry{2024.06--2024.08}{\textbf{Visiting Researcher}, Robotics and Systems Laboratory, ETH Zurich, Switzerland}
\cventry{2022.04--2025.01}{\textbf{Project Assistant Professor}, Department of Mechano-Informatics, Graduate School of Information Science and Technology, The University of Tokyo, Japan}

\section*{Education}
\cventry{2019.04--2022.03}{\textbf{Ph.D. in Mechano-Informatics}, Graduate School of Information Science and Technology, The University of Tokyo}
\cventry{2017.04--2019.03}{\textbf{M.S. in Mechano-Informatics}, Graduate School of Information Science and Technology, The University of Tokyo}
\cventry{2013.04--2017.03}{\textbf{B.S. in Mechano-Informatics}, Faculty of Engineering, The University of Tokyo}

\section*{Selected Publications}
\begin{itemize}
  \item K. Kawaharazuka et al., \textit{MEVIUS2: Practical Open-Source Quadruped Robot with Sheet Metal Welding and Multimodal Perception}, IEEE Robotics and Automation Practice (RAP), 2026.
  \item K. Kawaharazuka, K. Okada, and M. Inaba, \textit{Characteristics, Management, and Utilization of Muscles in Musculoskeletal Humanoids: Empirical Study on Kengoro and Musashi}, Advanced Intelligent Systems (AISY), 2026.
  \item K. Kawaharazuka et al., \textit{Vision-Language-Action Models for Robotics: A Review Towards Real-World Applications}, IEEE Access, 2025.
  \item K. Kawaharazuka, K. Okada, and M. Inaba, \textit{GeMuCo: Generalized Multisensory Correlational Model for Body Schema Learning}, IEEE Robotics and Automation Magazine (RAM), 2024.
  \item K. Kawaharazuka et al., \textit{Real-World Robot Applications of Foundation Models: A Review}, Advanced Robotics, 2024. (Advanced Robotics Best Survey Paper Award)
  \item K. Kawaharazuka et al., \textit{Toward Autonomous Driving by Musculoskeletal Humanoids: Study of Developed Hardware and Learning-Based Software}, IEEE Robotics and Automation Magazine (RAM), 2020.
\end{itemize}

\section*{Selected Grants/Funding}
\cventry{2024.10--2030.03}{\textbf{CRONOS} (Co-Investigator), Japan Science and Technology Agency (JST)}
\cventry{2024.10--2032.03}{\textbf{FOREST Program / Souhatsu} (PI), Japan Science and Technology Agency (JST)}
\cventry{2023.04--2027.03}{\textbf{JSPS KAKENHI Grant-in-Aid for Scientific Research (B)} (PI), Japan Society for the Promotion of Science (JSPS)}
\cventry{2023.04--2026.03}{\textbf{JSPS KAKENHI Challenging Exploratory Research} (PI), Japan Society for the Promotion of Science (JSPS)}
\cventry{2022.04--2024.03}{\textbf{JSPS KAKENHI Grant-in-Aid for Research Activity Start-up} (PI), Japan Society for the Promotion of Science (JSPS)}
\cventry{2020.12--2024.03}{\textbf{JST ACT-X and ACT-X Acceleration Phase} (PI), Japan Science and Technology Agency (JST)}
\cventry{2019.04--2022.03}{\textbf{JSPS Research Fellowship for Young Scientists (DC1)}, Japan Society for the Promotion of Science (JSPS)}

\section*{Selected Awards/Honors}
\cventry{2026.02}{\textbf{Funai Academic Award}, Funai Foundation for Information Technology}
\cventry{2025.11}{\textbf{Innovators Under 35 Japan 2025}, MIT Technology Review}
\cventry{2025.09}{\textbf{Advanced Robotics Best Survey Paper Award}, Advanced Robotics}
\cventry{2024.11}{\textbf{Mike Stillman Award}, IEEE-RAS International Conference on Humanoid Robots}
\cventry{2024.05}{\textbf{Best Conference Paper Award}, IEEE International Conference on Robotics and Automation (Open X-Embodiment Collaboration)}

\section*{Selected Professional Activities}
\cventry{2026}{Committee Member, \textbf{AI Robotics Strategy Council}, Ministry of Economy, Trade and Industry (METI), Japan}
\cventry{2026}{Special Track AI and Robotics Chairs, \textbf{IJCAI-ECAI 2026}}
\cventry{2025}{Organizer, \textbf{Workshop on Foundation Models for Robotic Design}, IEEE/RSJ IROS 2025}
\cventry{2025}{Organizer, \textbf{Dexterous Humanoid Manipulation Workshop}, IEEE-RAS Humanoids 2025}
\cventry{2025}{Organizer, \textbf{Workshop on Open Hardware in the Era of Robot Learning}, CoRL 2025}

\section*{Selected Open Source Projects}
\begin{itemize}
  \item \textbf{MEVIUS2} (2026): Open-source quadruped robot with sheet-metal construction and multimodal perception. \\
  \href{https://haraduka.github.io/mevius2-hardware/}{https://haraduka.github.io/mevius2-hardware/}
  \item \textbf{MEVITA} (2025): Open-source bipedal robot assembled from e-commerce components. \\
  \href{https://haraduka.github.io/mevita-hardware/}{https://haraduka.github.io/mevita-hardware/}
  \item \textbf{MEVIUS} (2025): Easily constructed open-source quadruped platform using sheet-metal welding and machining. \\
  \href{https://haraduka.github.io/mevius-hardware/}{https://haraduka.github.io/mevius-hardware/}
\end{itemize}

\section*{Publications}

\subsection*{Journal Articles (Peer Reviewed)}

% journal_replace_by_python

\subsection*{International Conference Proceedings (Peer Reviewed)}

% proceedings_replace_by_python

\end{document}
